\documentclass[10pt,a4paper]{amsart}
\usepackage[margin=19mm]{geometry}
\usepackage{amsmath,amssymb,mathtools}
\usepackage[scr=rsfs,cal=euler]{mathalfa}
\usepackage{xfrac}
\usepackage[unicode,hidelinks,bookmarks=false]{hyperref}
\newcommand{\ie}{i.e.\ }
\newcommand{\cf}{cf.\ }
\newcommand{\resp}{respectively\ }
\newcommand{\Subsection}{\subsection}
\providecommand{\nobreakdash}{\nobreak\hbox{-}}
\setlength{\emergencystretch}{2em}
\multlinegap0pt
\usepackage{bm}
\usepackage{bbm}
\usepackage{dsfont}
\usepackage{xspace}
\usepackage{cancel}
\usepackage{mathtools}
\usepackage{amsthm}
\makeatletter
\@ifundefined{lemma}{\newtheorem{lemma}{Lemma}}{}
\makeatother
\makeatletter
\expandafter\ifx\csname algbox\endcsname\relax
\newenvironment{algbox}[0]{\vskip 0.2in
\noindent 
\begin{fminipage}{6.3in}
}{
\end{fminipage}
\vskip 0.2in
}
\fi
\expandafter\ifx\csname fminipage\endcsname\relax
\newenvironment{fminipage}%
  {\begin{Sbox}\begin{minipage}}%
\fi
\expandafter\ifx\csname nscenter\endcsname\relax
\newenvironment{nscenter}
 {\parskip=0pt\par\nopagebreak\centering}
 {\par\noindent\ignorespacesafterend}
\fi
\renewcommand{\paragraph}{%
  \@startsection{paragraph}{4}%
  {\z@}{1.25ex \@plus 1ex \@minus .2ex}{-1em}%
  {\normalfont\normalsize\bfseries}%
}
\newcommand{\s}[1]{\mathsf{#1}}
\expandafter\ifx\csname solution\endcsname\relax
\newenvironment{solution}[1][Solution]{
    \begin{proof}[#1]
  }{
    \end{proof}
  }
\fi
\makeatother
\title[Recovery of Planted Subgraphs]{Recovery of Planted Subgraphs}
\author{Wasim Huleihel}
\date{Source: arXiv:2607.00897v1 (CC BY 4.0)}
\begin{document}
\maketitle

\noindent\emph{Source and licence.} Recovery of Planted Subgraphs. Version arXiv:2607.00897v1 (CC BY 4.0), distributed under the Creative Commons Attribution 4.0 International licence, as recorded in the arXiv licence metadata for this exact version and shown on the abstract page https://arxiv.org/abs/2607.00897v1. Full rights notice: licenses/arxiv-2607.00897-CC-BY-4.0.txt.\par\smallskip

\noindent\textit{Editorial scope.} Counted unit: the Lemma whose source label is \texttt{lem:anchor-avg-emb} in the article (source0.tex lines 2495--2506, source label \texttt{lem:anchor-avg-emb}), delivered together with the article's own complete original proof of it (source0.tex lines 2514--2628, one \texttt{proof} environment of 115 lines). No mathematics is written, repaired or re-derived: statement and proof are the article's own text line for line. The proof is detached from the statement in the source: the article states the result at lines 2495--2506 and prints its proof at lines 2514--2628, and the proof environment's own header names the statement, so the association is the article's own. Required context retained uncounted: none. Every definition, display and label of those lines is retained unchanged. Omitted from this package and not redistributed: 1--2494, 2507--2513, 2629--4429. Nothing inside a retained excerpt is omitted or altered: every retained body is delivered exactly as the article's lines are written, so no omission receipt of any kind accompanies this package. Citations: the retained excerpts carry no citation command at all, so no bibliography entry of the article is transcribed and no reference is redistributed. Reference adaptation: none was needed and none was made; every reference inside the delivered unit resolves inside it, so no unresolved reference or citation is produced. Printed slips kept literal: no printed word, symbol, hypothesis or number of the counted statement or of its proof was altered. Layout and class: the article's own class is \texttt{article} with packages including hyperref, color, graphicx, subfigure, amsmath, amssymb, amsfonts, bm, epsfig, epsf, url, dsfont, amsthm, mathrsfs; the wrapper is the lane's pinned \texttt{amsart} editorial wrapper at 10pt on A4, which restates the theorem-like environments the retained text needs and adds the article's own macro definitions that the retained text needs (\texttt{\textbackslash paragraph}, \texttt{\textbackslash s}), with extra wrapper packages (bm, bbm, dsfont, xspace, cancel, mathtools). The delivered numbering is the unit's own. Assets: no retained span contains a figure environment, a graphics-inclusion command, a PDF-page inclusion, a table or a TikZ picture; no image file is copied into this package, the package declares no asset dependency and no image file is redistributed with it. Licence: version arXiv:2607.00897v1 of this article is distributed under the Creative Commons Attribution 4.0 International licence, as recorded in the arXiv licence metadata for this exact version and shown on the abstract page https://arxiv.org/abs/2607.00897v1. A full rights notice is delivered at licenses/arxiv-2607.00897-CC-BY-4.0.txt. Characters: every retained excerpt is pure ASCII, so the delivered engine, XeLaTeX with its default Latin Modern text font, reports no missing character.
\begin{lemma}\label{lem:anchor-avg-emb}
Let $\Gamma_n$ be a simple graph on $k$ vertices with average degree
\begin{align}
\bar{d}(\Gamma_n)=\frac{1}{k}\sum_{u\in v(\Gamma_n)}\s{deg}_{\Gamma_n}(u)=\frac{2|e(\Gamma_n)|}{k}.
\end{align}
Fix integers $t\ge 2$ and $d\ge t-1$. There exists an absolute constant $C\ge 1$ such that
\begin{equation}\label{eq:anchor-avg-emb}
\frac{1}{k}  \s{Emb}_{t,d}(\Gamma_n)
\leq 
\big(Ct\big)^{  C(d+t)}  \bar{d}(\Gamma_n)^{  t-1}.
\end{equation}
\end{lemma}

\begin{proof}[Proof of Lemma~\ref{lem:anchor-avg-emb}] Let $\mathfrak H_{t,d}$ denote the family of connected rooted patterns on $t$ vertices and $d$ edges (root distinguished but otherwise labeled). The proof proceeds in five steps.

\paragraph{1) Exploration schemes.}
For a connected rooted pattern $H\in\mathfrak H_{t,d}$, fix a rooted spanning tree $T\subseteq H$ of size $t-1$, and fix an \emph{exploration order} $r=v_1,v_2,\dots,v_t$ in which each $v_{s+1}$ is adjacent in $T$ to some earlier vertex among $v_1,\dots,v_s$.
An \emph{exploration scheme} $\mathsf S$ comprises:
(i) the rooted spanning tree $T$;
(ii) the exploration order;
(iii) the choice of the extra (non-tree) edges of $H$ (there are $d-(t-1)$ of them).
Let $\mathcal S_{t,d}$ be the set of all exploration schemes over all $H\in\mathfrak H_{t,d}$.

A crude combinatorial bound suffices:
choose $T$ (at most $t^{  t-2}$ rooted labeled trees by Cayley),
choose an order (at most $t!$),
and choose the extra edges (at most $\binom{\binom{t}{2}}{d-(t-1)}\le (Ct)^{  2(d-(t-1))}$).
Thus
\begin{equation}\label{eq:scheme-count}
|\mathcal S_{t,d}|
\leq 
t^{  t-2}\cdot t!\cdot (Ct)^{  2(d-(t-1))}
\leq  (Ct)^{  C(d+t)}.
\end{equation}

\paragraph{2) Partial histories and the boundary recursion.}
Fix a relabeling $\sigma$ of $v(\Gamma_n)$, and write $\Gamma_n^\sigma$ for the relabeled graph.
For an exploration scheme $\mathsf S\in\mathcal S_{t,d}$ and an anchor $v\in v(\Gamma_n^\sigma)$, a \emph{partial history of length $s$} ($1\le s\le t$)
is a choice of an injective map sending $v_1,\dots,v_s$ to distinct vertices of $v(\Gamma_n^\sigma)$ that respects the (tree) adjacencies required by $\mathsf S$ among $v_1,\dots,v_s$.
Let $\mathcal F_s(\sigma)$ be the multiset of all partial histories of length $s$, formed by ranging over all anchors $v\in v(\Gamma_n^\sigma)$ and all schemes $\mathsf S\in\mathcal S_{t,d}$.

For $f\in\mathcal F_s(\sigma)$ write $U(f)\subseteq v(\Gamma_n^\sigma)$ for the current image (so $|U(f)|=s$), and
\begin{align}
\partial_{\Gamma_n^\sigma}(U)\triangleq\big|\{  \{x,y\}\in e(\Gamma_n^\sigma): x\in U,\ y\notin U  \}\big|
\end{align}
for the (undirected) edge boundary size of $U$.
Each extension from $s$ to $s+1$ chooses the next image vertex $v_{s+1}$ outside $U(f)$ that is adjacent, in $\Gamma_n^\sigma$, to the specified predecessor of $v_{s+1}$ in the tree $T$; consequently, for every $f\in\mathcal F_s(\sigma)$,
the number of admissible choices is at most $\partial_{\Gamma_n^\sigma}(U(f))$.
Summing over all $f$ yields the recursion
\begin{equation}\label{eq:histories-recursion-pointwise}
|\mathcal F_{s+1}(\sigma)|
\leq \sum_{f\in\mathcal F_s(\sigma)} \partial_{\Gamma_n^\sigma}\big(U(f)\big).
\end{equation}
Indeed, non-tree edges in $\mathsf S$ impose additional adjacency constraints and thus only \emph{decrease} the number of admissible choices, so \eqref{eq:histories-recursion-pointwise} remains valid.

\paragraph{3) Averaging over relabelings.}
By construction, $\s{Emb}_{t,d}(\Gamma_n)$ is invariant under relabelings of $v(\Gamma_n)$, so
\begin{align}
\s{Emb}_{t,d}(\Gamma_n)&=\mathbb E_\sigma\big[\s{Emb}_{t,d}(\Gamma_n^\sigma)\big]\\
&\leq \mathbb E_\sigma\big[  |\mathcal F_t(\sigma)|  \big],
\end{align}
because every full root-preserving embedding (for some $H$ and some $\mathsf S$) gives rise to at least one full history of length $t$.
Taking expectations of \eqref{eq:histories-recursion-pointwise}, we are reduced to bounding
\begin{align}
\mathbb E_\sigma\left[\sum_{f\in\mathcal F_s(\sigma)} \partial_{\Gamma_n^\sigma}\big(U(f)\big)\right]
\end{align}
in terms of $\mathbb E_\sigma\big[  |\mathcal F_s(\sigma)|  \big]$. We have the following lemma.

\begin{lemma}[Average boundary at size $s$]\label{cl:avg-boundary}
For each $1\le s\le t-1$,
\begin{equation}\label{eq:avg-boundary}
\mathbb E_\sigma\left[\sum_{f\in\mathcal F_s(\sigma)} \partial_{\Gamma_n^\sigma}\big(U(f)\big)\right]
\leq 
s  \bar{d}(\Gamma_n)\cdot \mathbb E_\sigma\big[  |\mathcal F_s(\sigma)|  \big].
\end{equation}
\end{lemma}

\begin{proof}[Proof of Lemma~\ref{cl:avg-boundary}]
For any $U\subseteq v(\Gamma_n^\sigma)$ we have $\partial_{\Gamma_n^\sigma}(U)\le\sum_{u\in U}\s{deg}_{\Gamma_n^\sigma}(u)$.
Thus
\begin{align}
\sum_{f\in\mathcal F_s(\sigma)} \partial_{\Gamma_n^\sigma}(U(f))
&\leq \sum_{f\in\mathcal F_s(\sigma)} \sum_{u\in U(f)} \s{deg}_{\Gamma_n^\sigma}(u)\\
&=\sum_{u\in v(\Gamma_n^\sigma)} \s{deg}_{\Gamma_n^\sigma}(u)\cdot N_s(u;\sigma),
\end{align}
where $N_s(u;\sigma)\triangleq|\{  f\in\mathcal F_s(\sigma): u\in U(f)  \}|$ is the multiplicity with which $u$ appears among the size-$s$ images.
By vertex-exchangeability under the uniform relabeling $\sigma$, $\mathbb E_\sigma[N_s(u;\sigma)]$ is the \emph{same} for all $u$.
Moreover,
\begin{align}
\sum_{u\in v(\Gamma_n^\sigma)} N_s(u;\sigma)&=\sum_{f\in\mathcal F_s(\sigma)} |U(f)|\\
&=s  |\mathcal F_s(\sigma)|,
\end{align}
hence $\mathbb E_\sigma[N_s(u;\sigma)]=(s/k)  \mathbb E_\sigma[|\mathcal F_s(\sigma)|]$ for every $u$. Taking expectations and using the fact that $\sum_{u}\s{deg}_{\Gamma_n^\sigma}(u)=2|e(\Gamma_n)|=k  \bar{d}(\Gamma_n)$ gives
\begin{align}
\mathbb E_\sigma\left[\sum_{f\in\mathcal F_s(\sigma)} \partial_{\Gamma_n^\sigma}(U(f))\right]
&\leq\sum_{u} \s{deg}_{\Gamma_n}(u)\cdot \frac{s}{k} \cdot \mathbb E_\sigma[  |\mathcal F_s(\sigma)|  ]\\
&=s  \bar{d}(\Gamma_n)\cdot\mathbb E_\sigma[  |\mathcal F_s(\sigma)|  ],
\end{align}
proving \eqref{eq:avg-boundary}.
\end{proof}

\paragraph{4) Induction and base size.}
Combining \eqref{eq:histories-recursion-pointwise} and Claim~\ref{cl:avg-boundary}, and iterating for $s=1,2,\dots,t-1$, we obtain
\begin{align}\label{eq:Ft-bound}
\mathbb E_\sigma\big[  |\mathcal F_t(\sigma)|  \big]
&\leq
\left(\prod_{s=1}^{t-1} s  \bar{d}(\Gamma_n)\right)\cdot \mathbb E_\sigma\big[  |\mathcal F_1(\sigma)|  \big]\\
&\leq t!  \bar{d}(\Gamma_n)^{  t-1}\cdot \mathbb E_\sigma\big[  |\mathcal F_1(\sigma)|  \big].
\end{align}
A partial history of length $1$ consists of choosing an anchor $v\in v(\Gamma_n^\sigma)$ and a scheme $\mathsf S\in\mathcal S_{t,d}$; hence $|\mathcal F_1(\sigma)|=k  |\mathcal S_{t,d}|$, deterministically.
Using \eqref{eq:scheme-count} and $t!\le (Ct)^t$, \eqref{eq:Ft-bound} yields
\begin{align}
\mathbb E_\sigma\big[  |\mathcal F_t(\sigma)|  \big]
\leq
k\cdot (Ct)^{  C(d+t)}\cdot \bar{d}(\Gamma_n)^{  t-1}.
\end{align}

\paragraph{5) From histories to embeddings.}
Every root-preserving embedding counted by $\s{Emb}_{t,d}(\Gamma_n^\sigma)$ corresponds to at least one full history (for some $\mathsf S$), so
$\s{Emb}_{t,d}(\Gamma_n^\sigma)\le |\mathcal F_t(\sigma)|$ and hence
\begin{align}
\s{Emb}_{t,d}(\Gamma_n)
&=\mathbb E_\sigma\big[\s{Emb}_{t,d}(\Gamma_n^\sigma)\big]\\
&\leq \mathbb E_\sigma\big[  |\mathcal F_t(\sigma)|  \big]\\
&\leq k\cdot (Ct)^{  C(d+t)}\cdot \bar{d}(\Gamma_n)^{  t-1}.
\end{align}
Dividing both sides by $k$ gives \eqref{eq:anchor-avg-emb}.
\end{proof}

\end{document}
